% SPDX-License-Identifier: Apache-2.0
% covers: Remote, Ask, Answer
\datasheet{Remote}{A memory-mapped bridge forwarding bus transactions over channels to an external software endpoint.}

\dsfacts{\code{lib/parts/src/remote.rs}}{\code{txhdl\_parts::remote::Remote}}%
{\code{//docs:examples}}

\subsection*{Function}

\code{Remote} connects to an AXI-Lite interconnect as a memory-mapped target
without decoding register state internally. Every transaction accepted by the
core is forwarded on \code{out} as an \code{Ask} message, and the corresponding
response returns on \code{back} as an \code{Answer}. The remote endpoint is
decoupled from the peripheral logic: in simulation, it connects to a software
model within the test process; on physical hardware, it connects via packet
frames to software running on a remote host.

This architecture enables hardware-software co-development. A target peripheral
can be developed and verified as software against the actual system bus. Once
the software implementation is proven, a dedicated hardware module replaces it
without modifying host bus transactions. Issue 297 documents this workflow.

The bus interface presents a multi-cycle peripheral latency: exactly one
transaction remains outstanding at a time, and the transaction channel is held
until the response arrives. This adheres to AXI-Lite protocol semantics and
accommodates network transit times.

\subsection*{Parameters}

\begin{center}\footnotesize
\begin{tabular}{@{}lp{0.7\textwidth}@{}}
\toprule
Parameter & Meaning \\
\midrule
\code{T} & The timeout limit in clock cycles before the peripheral responds
with \code{SLVERR}. This prevents an unresponsive remote endpoint from
permanently hanging the system bus. The datasheet tables use 40.
\code{ex\_remote} uses 1000 cycles, covering round-trip latency through two MACs.
\code{Board} sets this to 100,000,000 cycles (one second at 100~MHz), providing
sufficient margin for network round trips to a host workstation. \\
\bottomrule
\end{tabular}
\end{center}

\subsection*{Register map}

None. Every address within the peripheral's allocated aperture is interpreted
by the target program, and the target address is forwarded verbatim in the
\code{Ask} payload.

\dstables{Remote}

\subsection*{Behaviour}

\begin{itemize}
\item Exactly one transaction is outstanding at a time; \code{busy} indicates
active transfer, while \code{writing} indicates transaction direction.
\item Write transactions take priority over read transactions. Because writes
and reads arrive on separate channels, prioritizing writes prevents read streams
from starving pending store operations.
\item Each transaction includes a \code{tag} that increments with each dispatched
request. Responses arriving on \code{back} are accepted only when their tag matches
the currently outstanding transaction.
\item Tag validation guarantees safe timeout handling. When a transaction times
out after \code{T} cycles, the peripheral returns \code{SLVERR} to the bus and
retires the transaction. If a delayed response arrives later, its tag no longer
matches the active transaction, and the core discards it rather than misattributing
it to subsequent bus traffic.
\item The transaction tag advances when a transaction completes rather than when
it is issued, cleanly separating late responses from the subsequent request.
\item Read transactions complete on \code{bus\_r} with the returned data word,
and write transactions complete on \code{bus\_b}; either response returns
\code{SLVERR} if the remote target indicates failure or if the cycle timer expires.
\end{itemize}

\subsection*{Verification}

\begin{itemize}
\item Four unit tests in \code{lib/parts/src/remote.rs} (\code{//lib/parts:parts\_test}):
\begin{itemize}
\item \code{a\_word\_written\_to\_the\_program\_is\_read\_back\_from\_it}
verifies standard loopback;
\item \code{the\_program\_says\_which\_transactions\_fail} validates
remote error signaling;
\item \code{a\_slow\_program\_is\_waited\_for} tests latency tolerance;
\item \code{a\_program\_that\_never\_answers\_does\_not\_stop\_the\_bus}
confirms timeout recovery.
\end{itemize}
\item \code{//lib/examples:ex\_remote} executes the complete pipeline comprising
\code{RemoteLink} and dual Ethernet MACs. Netlists synthesized from this design
are simulated under both nvc and Verilator via
\code{//docs:remote\_sim\_remote\_remote\_tb\_test} and
\code{//docs:remote\_vsim\_remote\_test}.
\item \code{//cpu/vreteno:board\_test} executes the RISC-V core against the board
design: software writes a data word to \code{0x3300} and reads it back, with the
testbench harness serving response frames.
\end{itemize}

\subsection*{Used by}

\code{ex\_remote}, and \code{Board} at address \code{0x3300}, where a
\code{RemoteLink<1>} connects it to the board's Ethernet MAC for hardware-software
co-development over the network.

\subsection*{Limits}

The block supports AXI-Lite only with a single outstanding transaction, causing
bus accesses to block for the full remote round-trip latency. It implements no
transport layer directly: transmitting \code{Ask} requests to an external machine
requires \code{RemoteLink} and an Ethernet MAC, or equivalent channel adapters.
The timeout counter measures clock cycles rather than absolute wall-clock time;
operating at lower clock frequencies increases the elapsed timeout duration proportionally.
